% Fichier complété par tools/complete_missing_corrections.py.
% Source : DYN/DYN-04-TorseurDynamique/03_TT/03_TT.tex
% Statut : rédigé (3 question(s)).
\begin{corrigebox}[Corrigé rédigé]
\CorrigeQuestion{1}
$\{\mathcal C(1/0)\}=\{m_1\dot\lambda\vec i_0;\vec0\}$ et
                $\{\mathcal C(2/0)\}=\{m_2(\dot\lambda\vec i_0+
                \dot\mu\vec j_0);\vec0\}$, transportés au point demandé.
\CorrigeQuestion{2}
$\{\mathcal D(1/0)\}=\{m_1\ddot\lambda\vec i_0;\vec0\}$ et
                $\{\mathcal D(2/0)\}=\{m_2(\ddot\lambda\vec i_0+
                \ddot\mu\vec j_0);\vec0\}$.
\CorrigeQuestion{3}
Le torseur dynamique de l'ensemble est la somme : résultante
                $[(m_1+m_2)\ddot\lambda]\vec i_0+m_2\ddot\mu\vec j_0$ et moment égal à
                la somme des moments transportés en $B$.
\end{corrigebox}
